% =========================================================
\section{Objectifs de l'automatisation}

\subsection{Introduction}

Automatiser une machine ou une installation n'est pas une fin en soi. Une entreprise automatise pour répondre à des besoins précis : améliorer la production, garantir la qualité, réduire les coûts, protéger les personnes ou surveiller un procédé.

L'automatisation doit donc être justifiée par un ou plusieurs objectifs clairement définis dans le cahier des charges.

\subsection{Pourquoi automatiser ?}

On automatise principalement pour :

\begin{itemize}
    \item effectuer des opérations répétitives sans fatigue ;
    \item obtenir une production régulière et conforme ;
    \item augmenter la cadence de production ;
    \item réduire les erreurs humaines ;
    \item surveiller en permanence certaines grandeurs ;
    \item protéger les opérateurs contre les tâches dangereuses ;
    \item commander plusieurs équipements selon une séquence précise ;
    \item réduire certaines consommations d'énergie ou de matières.
\end{itemize}

\subsection{Objectifs techniques}

Sur le plan technique, l'automatisation permet :

\begin{itemize}
    \item la répétitivité des opérations ;
    \item la précision des actions ;
    \item la rapidité d'exécution ;
    \item la coordination de plusieurs actionneurs ;
    \item la surveillance continue du système ;
    \item la gestion de séquences complexes.
\end{itemize}

\paragraph{Exemple :} une ligne de remplissage peut remplir, peser, boucher et étiqueter des bouteilles dans un ordre précis, sans intervention permanente de l'opérateur.

\subsection{Objectifs économiques}

L'automatisation cherche souvent à améliorer la rentabilité de l'installation.

Elle peut permettre de :

\begin{itemize}
    \item réduire le temps de cycle ;
    \item diminuer les rebuts et les pertes de matière ;
    \item réduire certaines interventions manuelles ;
    \item augmenter le rendement des machines ;
    \item améliorer la disponibilité de l'installation ;
    \item produire avec une qualité plus constante.
\end{itemize}

Cependant, l'automatisation représente aussi un investissement : achat du matériel, câblage, programmation, formation du personnel et maintenance.

\subsection{Objectifs liés à la sécurité et aux conditions de travail}

Un système automatisé bien conçu peut :

\begin{itemize}
    \item éloigner l'opérateur des zones dangereuses ;
    \item supprimer certaines tâches pénibles ou répétitives ;
    \item détecter plus rapidement un défaut ou une anomalie ;
    \item déclencher un arrêt d'urgence ou une alarme ;
    \item imposer une séquence de fonctionnement plus sûre.
\end{itemize}

L'opérateur devient alors davantage un superviseur qu'un exécutant direct.

\paragraph{Attention :} automatiser ne signifie pas supprimer toute surveillance humaine. La sécurité doit être conçue dès le départ et entretenue pendant toute la durée de vie de l'installation.

\subsection{Objectifs de qualité et de traçabilité}

L'automatisation permet d'obtenir :

\begin{itemize}
    \item une qualité régulière ;
    \item le respect strict d'un cahier des charges ;
    \item un contrôle automatique de certaines caractéristiques ;
    \item l'enregistrement de données de production ;
    \item la génération d'alarmes en cas de dérive.
\end{itemize}

Dans certaines industries, il est important de savoir quand, comment et avec quels paramètres un produit a été fabriqué. C'est le rôle de la traçabilité.

\subsection{Objectifs de flexibilité}

Un automatisme programmé, comme celui réalisé avec un API, peut être modifié plus facilement qu'un automatisme entièrement câblé.

Il devient alors possible de :

\begin{itemize}
    \item changer un format de production ;
    \item modifier une recette ou une séquence ;
    \item adapter la cadence ;
    \item ajouter une nouvelle étape ;
    \item passer d'un mode manuel à un mode automatique.
\end{itemize}

Cette flexibilité est très importante dans les industries où les produits ou les demandes changent souvent.

\subsection{Exemples industriels}

\begin{itemize}
    \item \textbf{Traitement de l'eau :} démarrage et arrêt automatiques des pompes selon le niveau des réservoirs.
    \item \textbf{Convoyeur :} détection des colis et commande du moteur.
    \item \textbf{Remplissage de cuve :} ouverture d'une électrovanne jusqu'à un niveau ou un volume déterminé.
    \item \textbf{Porte automatique :} détection d'une personne et commande du moteur d'ouverture.
    \item \textbf{Tri de pièces :} identification et orientation des pièces selon leur type.
    \item \textbf{Ascenseur simplifié :} gestion des appels, des niveaux et des sécurités.
\end{itemize}

\subsection{Avantages et limites de l'automatisation}

\begin{table}[h]
\centering
\begin{tabularx}{\textwidth}{|X|X|}
\hline
\textbf{Avantages} & \textbf{Limites / points de vigilance} \\
\hline
Augmentation de la productivité & Coût d'investissement initial \\
\hline
Qualité plus régulière & Complexité de conception et de maintenance \\
\hline
Réduction de certaines tâches pénibles & Besoin de personnel qualifié \\
\hline
Meilleure sécurité lorsqu'elle est bien conçue & Dépendance à l'énergie et à l'électronique \\
\hline
Surveillance continue du procédé & Risque d'obsolescence du matériel \\
\hline
Flexibilité grâce à la programmation & Nécessité de formation et de documentation \\
\hline
\end{tabularx}
\caption{Avantages et limites de l'automatisation}
\end{table}

\subsection{Conclusion}

Automatiser consiste à confier à un système technique des opérations qui étaient auparavant réalisées manuellement ou avec une aide limitée.

Les objectifs peuvent être techniques, économiques, humains ou liés à la qualité. Ils doivent être clairement définis avant de choisir les capteurs, les actionneurs, l'automate et le programme.

\subsubsection*{À retenir}

\begin{quote}
L'automatisation vise à améliorer la production, la qualité, la sécurité et la régularité d'un procédé, tout en tenant compte des coûts, de la maintenance et des compétences nécessaires.
\end{quote}

\paragraph{Schéma fondamental :}

\[
\text{Besoin} \rightarrow \text{Objectifs} \rightarrow \text{Choix techniques} \rightarrow \text{Réalisation}
\]